\section{Introduction}\label{sec:intro}
Each reduced rational $x=a/b$ has an occurrence set $D(x)=\{d\ge2: d^2\mid a,\ d\equiv1 \pmod b\}$. Writing the divisors in exponent coordinates over the primes whose squares divide $a$ identifies $D(x)$ with the nonzero points of a box intersected with a lattice $\Lambda_b$ of exponent vectors defined by one congruence (\cref{prop:fibre}). This finite set, with its Graver moves, parallelotopes and their union, is the geometry studied in the GEO-I releases of the UNNS programme~\cite{GEOI}: its Graver moves, the parallelotopes they span, and the union of those parallelotopes.

The question of this manuscript is what happens when the arithmetic moves. Under the collective map $x\mapsto\Cmap(x)=x+L^2/G$, with $L=\operatorname{lcm}D(x)$ and $G=\gcd\{d-1:d\in D(x)\}$ (the rational dynamics FIELD-I of the programme~\cite{FoundII}), the occurrence set can only grow, $D(x)\subseteq D(\Cmap(x))$, but the lattice and its Graver set change, and geometric structure that existed before need not exist afterwards, nor conversely. We call the study of such comparisons \emph{transport}. The manuscript treats it for explicitly defined source and target lattices and a fixed box; the evolution of whole orbits is the subject of a later manuscript.

\paragraph{Three objects.} For a lattice $\Lambda$ and box $Q$ we use three objects that are \emph{not} interchangeable: the abstract complex $\Kabs(\Lambda,Q)$ of cells, the realised union $U(\Lambda,Q)\subset\R^k$ of the corresponding closed parallelotopes, and, when the cells meet only in common faces, the embedded subdivision $\Kemb$. Every statement below names the object it concerns.

\paragraph{Results.}
\begin{enumerate}[label=(\Roman*)]
\item \emph{Exact criteria (Sections~\ref{sec:lattice}--\ref{sec:cells}).} Inclusion of relation lattices is decided from finitely many group data (\cref{thm:P2T1}). A cell of the source survives in the abstract target complex exactly when its vertices lie in the target fibre and its generators are Graver moves of the target (\cref{thm:cell-transport}); hence no cell of dimension larger than the rank of the common Graver moves survives. For the three-stage orbit $P^4/3\to2P^4/3\to P^4$ all positive-dimensional cells are lost for every number $k\ge2$ of primes (\cref{cor:threestage}), although the occurrence set strictly grows.
\item \emph{Abstract loss is not geometric loss (\cref{sec:realised}).} Every abstractly lost edge of that orbit still lies in the target union (\cref{thm:E28}), while some two-dimensional cells leave it (\cref{prop:E29}); four plausible monotonicity laws are refuted by exact witnesses.
\item \emph{Rectangular targets and the corner cube (\cref{sec:corner}).} For a target lattice $\bigoplus o_i\Z$ the uncovered region is exactly an open corner cube $O$; edges never meet it in any dimension, so the union is preserved for $k=2$ (\cref{thm:corner}), while for $k=3$ the failure is governed by two-dimensional cells and, for the even-sum source, occurs exactly when the scaled box has a side $\ge2$ (\cref{thm:evensum}).
\item \emph{Planar monotonicity (\cref{sec:planar}).} For every inclusion $\Lambda\subseteq\Lambda'$ of full-rank lattices of $\Z^2$ and every box, $U(\Lambda,Q)\subseteq U(\Lambda',Q)$ (\cref{thm:planar}). The proof isolates \emph{hard} cells, those that touch every coordinate hyperplane, and shows that in the plane their geometry forces coverage.
\item \emph{Dimensional obstruction (\cref{sec:obstruction}).} Containment fails in dimension three, in infinite families closed under diagonal dilation, with a target having no free coordinate (\cref{thm:C1}); and also for a target, the body-centred lattice, that is not split in any coordinate (\cref{prop:C3}, finite-exact).
\item \emph{Slab criterion (\cref{sec:slab}).} For a target that splits along one coordinate with period $k$, the union is an explicit stack of planar unions, and containment of a cell is a planar condition. The criterion as first published (closed height) is false; we prove and use the corrected form (\cref{thm:slab-ii}).
\end{enumerate}

\paragraph{What is not claimed.} Statements checked only by finite exact computation are labelled FINITE-EXACT and carry their domain. Nothing is asserted about arbitrary lattices outside the stated families, and no planar theorem is extended to dimension three. Several results are elementary; no priority claim is made. The classification of non-split three-dimensional targets that preserve the union is open (\cref{sec:open}).

\paragraph{Relation to earlier work.} The relation lattices, Graver cells and planar stabilisation theory are inherited from GEO-I~\cite{GEOI} and the Foundations volumes~\cite{FoundI,FoundII,RefColl} and are cited with their hypotheses; the planar stabilisation theorem G12-T8 is used only as a cited prerequisite (\cref{sec:g12}). Throughout, \texttt{CW-nnnn} denotes the row of the consolidated theorem crosswalk and the second identifier is the source theorem name; \cref{app:map} lists all of them.

\paragraph{Reading guide to the figures.} The geometric arguments are illustrated where they are made. \Cref{fig:objects} fixes the three objects; \cref{fig:equalfibres,fig:threestage,fig:planar} illustrate the transport of lattices and of abstract cells; \cref{fig:dim,fig:corner} the corner cube; \cref{fig:lost} the difference between abstract loss and containment; \cref{fig:planarproof} the planar theorem; \cref{fig:noncontain,fig:dilation,fig:bcc} the three-dimensional obstructions; \cref{fig:slab,fig:strict} the slab criterion; \cref{fig:hardcell,fig:f110} hard cells and the coupled specimen. Each caption states whether the picture is an \emph{exact computation} (coordinates produced and checked in exact rational arithmetic by \texttt{scripts/figs\_*.py}; the data are in \texttt{results/FIGURE\_DATA.json}), a \emph{theorem-level construction} (exact coordinates of the construction in a proof) or a \emph{schematic}. Colours: source structures blue, targets teal, surviving relations green, lost relations red, uncovered regions light orange, witnesses magenta stars; every distinction is also carried by line style, hatching or marker shape.
